\chapter{研究设计}

本章在前文理论基础上进入具体研究设计，主要说明本文如何把 AI 文旅短视频的内容特征转化为可分析变量。章节依次提出研究模型与研究假设，交代数据来源、样本筛选、变量编码、模型设定和分析方法，为后续实证检验提供基础。

\section{研究模型与研究假设}

本文的研究目标是识别 AI 文旅短视频内容特征对传播效果的影响。根据第二章的理论分析，传播效果并不只取决于是否使用 AI，而更取决于 AI 内容最终呈现出的技术质感、视觉风格以及视频本身的观看成本。基于 S-O-R 模型，本文将 AI 介入特征和视频内容特征视为外部刺激，将用户对内容的新奇感、审美感、真实感和互动意愿视为潜在心理机制，并以点赞、评论、收藏和转发等互动行为表征传播反应。

本文的解释变量由技术使用、内容呈现和传播条件三部分构成。技术使用部分关注 AI 在声音、画面和整体生成中的参与方式；内容呈现部分关注画面质量、主题取向、主体对象和风格表达；传播条件部分则关注视频时长与互动累积时间。通过这样的设置，本文试图避免用“是否使用 AI”过度概括复杂内容差异，并更细致地分析用户真正感知到的传播刺激。

在提出研究假设之前，需要先说明主要变量的编码依据。本文不把 AI 文旅短视频简单处理为“使用 AI”或“未使用 AI”，而是按照用户在视频中能够直接识别的内容刺激进行分类。AI 介入形式依据 AI 主要作用的模态分为 AI 音频、AI 画面和混合三类；AI 介入程度依据 AI 在视频中的参与范围和是否承担核心呈现功能分为浅度、中度和深度；AI 技术质感依据画面清晰度、运动自然度、人物或景观是否变形、生成痕迹是否影响观看等标准分为粗糙、普通和精良；视频风格依据整体视觉表达分为纪实/写实、卡通/艺术化和抽象/超现实；内容导向则依据视频主要传播目的分为信息实用型和娱乐审美型。上述编码既参考生成视频质量评价中关于视觉完整度、协调性和一致性的讨论，也结合文旅短视频中目的地形象呈现、地方符号表达和用户互动行为的特点。

结合既有文献和研究问题，本文提出如下研究假设。

\textbf{H1：AI 技术质感越高，传播效果越好。} 本文将 AI 技术质感编码为粗糙、普通和精良三个等级，主要判断依据是画面完成度、清晰度、自然度和生成痕迹是否影响观看。该假设并非来自研究者主观判断，而是建立在生成视频质量评价和 AIGC 营销研究基础上。生成视频质量研究指出，视觉协调性、视频文本一致性和领域适配会影响用户对生成视频的质量判断\cite{lit03}；AIGC 营销内容研究也表明，生成内容的价值感、可信度和情感连接会进一步影响用户参与\cite{lit18}。在文旅场景中，粗糙的 AI 画面容易造成违和感，而更高质量的画面更可能提升目的地形象呈现和用户互动意愿。

\textbf{H2：相较于纪实/写实风格，抽象/超现实风格更可能提升传播效果。} 本文将视频风格编码为纪实/写实、卡通/艺术化和抽象/超现实三类，其中抽象/超现实主要指现实拍摄难以直接呈现的想象化、奇观化或非现实表达。该假设主要来自旅游短视频视觉营销研究。既有研究表明，旅游短视频的视觉内容、视觉视角和叙事吸引力会影响旅游意向\cite{lit06}，TikTok 旅游短视频中的娱乐性和互动性也会通过沉浸体验影响用户行为意向\cite{lit05}。AI 技术能够生成现实拍摄难以实现的想象场景，因此抽象/超现实风格更可能形成视觉冲击和互动动机。

\textbf{H3：AI 介入程度会影响传播效果，但更深的 AI 介入未必一定带来更高互动。} 本文将 AI 介入程度编码为浅度、中度和深度，区分的是 AI 在视频中是局部辅助、重要参与还是整体生成。该假设来自 AI 内容可信度和旅游 AI 应用研究。AI 能够提升内容生产能力和个性化体验，但也可能带来可信度、真实感和目的地形象风险\cite{lit23}；用户对 AI 内容的接受也会受到感知可信度影响\cite{lit14}。AIGC 营销研究进一步表明，内容价值与可信度会共同影响用户参与\cite{lit18}。因此，AI 介入越深，视频越可能呈现更明显的生成特征，但用户是否互动还取决于最终质量、可信度和真实感。

\textbf{H4：娱乐审美型内容相较于信息实用型内容更可能获得较高传播效果。} 本文将内容导向编码为信息实用型和娱乐审美型，前者侧重景点介绍、活动信息、路线提示等功能性内容，后者侧重视觉审美、情绪表达和趣味传播。该假设来自旅游短视频内容类型和感知趣味研究。Wu 和 Ding 基于顾客启发理论比较了信息导向与情感导向旅游短视频，发现旅游短视频能够通过激发顾客启发影响旅游意向，且享乐型动机下的作用更强\cite{lit22}；旅游短视频内容营销研究也表明，内容价值会通过情感反应和目的地态度影响用户行为意向\cite{lit04}；主题公园 TikTok 研究进一步指出，短视频内容和感知趣味会影响访问意向\cite{lit17}。因此，娱乐审美型内容更可能通过视觉新奇感、情绪唤起和趣味体验带来互动。

\textbf{H5：视频时长会负向影响传播效果。} 本文将视频时长作为连续变量处理，并在模型中进行对数化转换。该假设来自短视频平台注意力和沉浸体验研究。短视频平台以短时、连续推荐和快速切换为重要使用场景，用户观看通常处于有限注意力分配之中\cite{lit21}；TikTok 旅游短视频研究也表明，娱乐性、互动性和沉浸体验会影响用户行为意向\cite{lit05}。因此，在同等内容质量下，过长视频会提高观看成本，可能降低完整观看和进一步互动的概率。需要说明的是，视频时长也可能存在“过短信息不足、适中效果较好、过长观看成本过高”的非线性关系，但本文样本的主回归和稳健性结果更支持负向线性判断，因此本文将 U 型或倒 U 型关系作为后续扩展检验方向，而不作为主假设。

上述假设并不意味着所有变量都应在回归中显著。本文更关注的是，在控制其他变量后，哪些因素仍然具有稳定的净影响。

\section{数据来源与样本筛选}

本文数据来自抖音平台省级文旅账号发布的短视频。研究以 33 个省份/地区文旅账号中筛选出的 AI 文旅短视频为基础，最终形成 722 条深度分析样本。由于不同省份原始采集表的采集深度、重复情况和字段完整度存在差异，本文不直接使用原始汇总表开展模型分析，而是在统一筛选、去重与复核后保留最终有效样本。论文中的正式统计分析、回归结果和机器学习补充分析均基于同一份 722 条正式分析样本。

纳入最终分析样本的视频同时满足两个条件：一是内容与旅游目的地、地方文化、城市形象、景区资源或文旅活动相关；二是在声音、画面或整体视听呈现中存在可识别的 AI 生成或 AI 辅助生成痕迹。经去重、筛选与复核后，最终形成 722 条深度分析样本，且最终样本链接无重复。

为提高样本筛选过程的可复现性，本文进一步将筛选规则细化为以下七个步骤。第一步，来源汇总。将 33 个省份/地区的原始采集表统一整理为候选池，并保留标题、链接、互动数据、发布时间和来源省份等基础字段。第二步，链接去重。以视频链接为主键识别重复记录，对于同一链接在不同表中重复出现的情况，仅保留一条记录，避免同一视频被重复计入样本。第三步，显性 AI 线索初筛。依据标题、话题词和描述中出现的“AI”“人工智能”“AIGC”“数字人”“AI 绘画”“AI 视频”等显性线索识别第一轮 AI 相关候选视频。第四步，文旅属性筛选。对初筛结果进行复核，仅保留与地方旅游资源、城市形象、景区景点、地方文化、节庆活动和文旅宣传直接相关的视频，删除虽含 AI 表述但与文旅主题无关的内容。第五步，AI 痕迹核验。对保留样本进一步判断其在声音、画面或整体视听结构中是否存在可识别的 AI 介入痕迹，例如 AI 配音、AI 绘图、AI 生成场景、虚拟人物或明显的生成式重构效果。第六步，异常样本清理。剔除链接失效、内容无法识别、同一视频多版本重复上传、标题与实际内容严重不符以及无法稳定判断是否属于研究对象的样本。第七步，最终复核。对进入深度分析阶段的样本再次核对发布时间、互动量字段和编码基础信息，形成最终 722 条有效样本。

上述规则的核心目的，是把“是否属于 AI 文旅短视频”这一判断拆分为“文旅属性成立”和“AI 痕迹成立”两个独立条件，并在此基础上进行去重、清理和复核。这样可以减少单纯依赖关键词筛选带来的误差，也有助于回应数据筛选标准不清晰的问题。

\begin{figure}[!htbp]
  \centering
  \vspace{-0.5em}
  \includegraphics[width=0.72\textwidth]{../logs/extended_work/06_数据筛选流程图.png}
  \vspace{-0.8em}
  \caption{样本筛选与复核流程}
  \label{fig:screening}
  \vspace{-0.6em}
\end{figure}

省级主样本的互动数据统一截取于 2026 年 4 月 13 日。考虑到点赞、评论、收藏和转发等互动指标均为累计值，采集时点会直接影响因变量大小，因此本文将 2026 年 4 月 13 日作为省级主样本的统一数据截面，并据此计算发布距采集日天数。最终深度分析样本的发布时间最早为 2023 年 4 月 12 日，最晚为 2026 年 4 月 13 日，覆盖 32 个省份/地区。

最终分析样本保留的字段包括省份、标题、链接、点赞量、评论量、收藏量、转发量、视频时长、视频发布日期以及人工编码得到的 AI 介入与视频内容变量。由于本文关注平台公开互动效果，未使用账号后台曝光量、完播率、投流信息等非公开数据。

\section{变量定义与编码}

本文变量体系包括因变量、核心解释变量和控制变量。因变量用于衡量传播效果，解释变量围绕 AI 使用、内容呈现和平台传播条件展开。

\subsection{因变量}

本文将传播效果定义为用户在短视频平台上的公开互动行为，并以点赞量、评论量、收藏量和转发量之和构造总互动量，即 $\text{总互动量}=\text{点赞量}+\text{评论量}+\text{收藏量}+\text{转发量}$。

由于总互动量存在明显右偏分布，本文主模型使用 $\log(\text{总互动量}+1)$ 作为因变量。加 1 是为了避免互动量为 0 时无法取对数。为检验结论稳健性，本文还分别以 $\log(\text{点赞量}+1)$、$\log(\text{评论量}+1)$、$\log(\text{收藏量}+1)$ 和 $\log(\text{转发量}+1)$ 作为替代因变量；同时构造日均互动量，即总互动量除以发布距采集日天数，并在稳健性检验中用于辅助判断采集时点差异是否改变核心结论。发布距采集日天数依据视频发布日期与统一采集截面之间的间隔计算，在建模处理中至少保留 1 天，以避免时间变量取值为 0。

\subsection{解释变量}

各变量的具体编码见表 \ref{tab:variables-ai} 和表 \ref{tab:variables-content}。表中列出的分类并不是为了机械扩大变量数量，而是为了把样本中不同层面的内容差异区分开来：有些变量描述生产方式，有些变量描述用户直接看到的画面与风格，还有一些变量描述文旅内容本身。

\begin{table}[!htbp]
  \centering
  \vspace{-0.5em}
  \caption{AI 相关变量定义与编码}
  \label{tab:variables-ai}
  \scriptsize
  \setlength{\tabcolsep}{3pt}
  \renewcommand{\arraystretch}{0.86}
  \begin{tabular}{L{0.20\textwidth}L{0.25\textwidth}L{0.43\textwidth}}
    \toprule
    变量 & 类别或形式 & 含义说明 \\
    \midrule
    AI 介入形式 & AI 音频、AI 画面、混合 & AI 主要作用于声音、画面或多种模态共同参与。 \\
    AI 介入程度 & 浅度、中度、深度 & AI 在视频生产和呈现中的参与深浅。 \\
    AI 技术质感 & 粗糙、普通、精良 & AI 生成内容的画面完成度、真实感和视觉质量。 \\
    \bottomrule
  \end{tabular}
  \vspace{-0.8em}
\end{table}

\begin{table}[!htbp]
  \centering
  \vspace{-0.5em}
  \caption{视频内容与控制变量定义}
  \label{tab:variables-content}
  \scriptsize
  \setlength{\tabcolsep}{3pt}
  \renewcommand{\arraystretch}{0.86}
  \begin{tabular}{L{0.20\textwidth}L{0.25\textwidth}L{0.43\textwidth}}
    \toprule
    变量 & 类别或形式 & 含义说明 \\
    \midrule
    视频主体 & 人、景色、物 & 视频中主要被呈现的对象。 \\
    视频主题 & 自然风光、历史人文、现代城市/生活、家国情怀、剧情/整活 & 视频内容的主要主题方向。 \\
    内容导向 & 信息实用型、娱乐审美型 & 内容更偏旅游信息服务还是审美娱乐传播。 \\
    视频风格 & 纪实/写实、卡通/艺术化、抽象/超现实 & 视频整体视觉和叙事风格。 \\
    视频时长 & 连续变量 & 视频持续时间，单位为秒，模型中取 $\log(\text{视频时长}+1)$。 \\
    发布距采集日天数 & 连续变量 & 视频发布日与统一采集日之间的间隔天数，模型中取 $\log(\text{发布距采集日天数}+1)$。 \\
    \bottomrule
  \end{tabular}
  \vspace{-0.8em}
\end{table}

AI 介入形式主要用于判断 AI 技术作用于视频的哪个环节。若视频主要使用 AI 配音、AI 音色、智能朗读或明显的合成声音，而画面仍以实拍、普通剪辑或常规图像为主，则编码为 AI 音频；若视频主要体现为 AI 绘图、AI 生成场景、AI 特效画面、虚拟人物或画面重构，而声音部分并非主要 AI 特征，则编码为 AI 画面；若声音和画面均存在较明显的 AI 参与，或者视频整体由多模态 AI 共同生成和包装，则编码为混合。该变量关注的是 AI 介入的主要模态，而不是判断视频质量高低。

AI 介入程度用于判断 AI 在视频生产中的作用深浅。浅度介入指 AI 主要承担局部辅助功能，例如个别配音、局部画面修饰、简单字幕或局部特效，视频主体仍主要依赖实拍或常规剪辑；中度介入指 AI 已经参与到较多画面、声音或叙事表达中，但视频仍保留较多真实素材或人工剪辑痕迹；深度介入指视频主要视觉效果、人物形象、场景构建或整体视听结构明显依赖 AI 生成，用户观看时能够直接感知到较强的生成式特征。该变量反映的是生产过程和呈现结构中的 AI 参与程度。

AI 技术质感是本文最重要的解释变量之一，用于判断 AI 生成内容最终呈现出来的质量。粗糙质感指画面存在明显失真、人物或景观变形、边缘模糊、动作不自然、清晰度较低或生成痕迹破坏观看体验；普通质感指画面基本完整，主要内容能够被正常识别，生成痕迹虽存在但不严重影响观看，整体视觉效果处于可接受水平；精良质感指画面完成度较高，人物、景观和运动变化较为自然，色彩、构图和视觉冲击力较强，AI 痕迹能够被转化为审美效果或新奇体验。该变量不以 AI 使用深浅为依据，而以用户最终看到的画面质量为依据。

视频主体用于判断视频画面中最主要的呈现对象。若内容主要围绕人物、虚拟人、游客或拟人化角色展开，则编码为人；若主要呈现自然景观、城市地标、街区空间、景区场景等环境对象，则编码为景色；若视频重点呈现文创产品、地方物产、非遗器物、建筑局部或其他具体物件，则编码为物。该变量用于区分文旅内容中被突出展示的对象类型。

视频主题用于判断视频表达的主要内容方向。自然风光主要包括山水、湖海、森林、季节景观等自然资源；历史人文主要包括历史遗迹、传统文化、非遗、古城古镇和地方文化符号；现代城市/生活主要包括城市地标、夜景、街区生活、消费场景和现代公共空间；家国情怀主要包括国家叙事、纪念活动、重大节庆和宏观情感表达；剧情/整活则主要指借助剧情反转、拟人化、网络梗或搞笑表达来完成传播。若一个视频同时包含多种主题，本文按照其最主要的叙事重点进行归类。

内容导向用于区分视频的传播目的。信息实用型主要以告知、介绍和服务为目标，例如景点介绍、活动预告、交通路线、游玩攻略、政策信息和旅游提示；娱乐审美型主要以吸引注意、唤起情绪和形成审美体验为目标，例如视觉奇观、情绪短片、城市氛围展示、拟人化表达和趣味化传播。该变量关注的是内容功能，而不是视频是否具有审美画面，因为信息实用型视频也可能使用精良画面。

视频风格用于判断视频整体视觉和叙事表达方式。纪实/写实风格指画面接近真实拍摄或真实场景再现，强调目的地的现实面貌；卡通/艺术化风格指视频采用插画、动画、漫画、国风绘画、艺术滤镜等明显艺术化表达，但仍保留较清晰的具象对象；抽象/超现实风格则指视频呈现出现实拍摄难以实现的想象场景、夸张变形、奇观化组合或强烈非现实感，例如地标拟人化、景观漂浮、跨时空组合和梦幻化场景。该变量主要判断用户看到的风格效果，而不是 AI 技术本身。

视频时长和发布距采集日天数属于连续变量。视频时长以秒为单位，用于衡量观看成本，模型中取 $\log(\text{视频时长}+1)$；发布距采集日天数依据视频发布日与统一采集日之间的间隔计算，用于控制公开互动量的时间累积效应，模型中取 $\log(\text{发布距采集日天数}+1)$。这两个变量不属于人工主观分类，但会影响用户互动和累计互动规模，因此纳入模型控制。

\subsection{变量编码可靠性控制}

由于 AI 介入程度、AI 技术质感、视频风格和内容导向等变量带有一定内容判断属性，本文在编码过程中设置了统一的判定口径。编码前先依据变量定义形成编码说明，明确各类别的边界特征；编码时同时参考视频标题、画面内容、声音形式、AI 痕迹和文旅主题关联，而不以单一关键词作为判断依据。对于介于两个类别之间的样本，本文按照更保守的原则归类，并在样本复核阶段统一处理相似情形。

为降低人工编码的随意性，本文对完成编码后的样本进行了复核。复核重点包括 AI 介入是否可识别、视频是否具备文旅属性、技术质感分类是否与画面完成度相符，以及娱乐审美型与信息实用型之间是否存在混淆。对于无法稳定判断为研究对象的样本，不纳入最终有效样本。考虑到本研究主要基于单人编码完成，且 AI 辅助复核在高主观性变量上仍可能与人工判断存在分歧，本文未将其表述为严格意义上的双人一致性检验结果。现阶段的做法是，在不查看互动结果的条件下，使用统一编码说明对部分样本进行辅助复核，用于识别分歧样本、修订变量边界并检查高主观性变量的稳定性；最终正式分析仍以统一规则下的人工作业结果为准。本文主要通过统一编码说明、集中复核、疑难样本保守归类和最终样本清理来尽量降低主观判断误差，并在研究不足中保留这一方法限制。

此外，本文在编码操作中将因变量互动数据与人工内容变量编码分离处理。人工编码阶段使用不显示点赞量、评论量、收藏量、转发量和总互动量等结果指标的编码表，编码者主要依据视频内容本身、AI 介入痕迹、视觉质感、主题与风格进行判断；互动数据则在内容编码完成后再与编码结果合并用于统计分析。通过这种结果变量隔离的方式，可以减少高互动或低互动结果对人工分类判断的暗示作用，降低编码过程中的主观预期偏差。为便于复现，本文建议将数据清洗脚本、变量编码说明、模型估计脚本和不含敏感信息的脱敏样本字段上传至 GitHub 仓库，并在论文附录或脚注中提供仓库地址。公开内容应排除 cookies、浏览器会话、请求头等私密信息，原始研究数据若不便完全公开，可至少公开字段说明、筛选流程、示例数据和可复现实验代码。

\subsection{参照组设置}

分类变量进入回归模型时需要转化为虚拟变量。本文以各变量中较基础或更常见的类别作为比较基准：声音层面的 AI 使用、浅度介入、粗糙质感、人作为主体、自然风光主题、信息实用型内容和纪实/写实风格分别作为参照组。连续变量视频时长和发布距采集日天数不设参照组。

\section{模型设定}

本文主模型为多元 OLS 回归模型，形式为
\begin{compactformula}
\begin{equation}
\begin{aligned}
Y_i={}&\beta_0+\beta_1 AIForm_i+\beta_2 AIDegree_i+\beta_3 AIQuality_i
+\beta_4 Subject_i+\beta_5 Theme_i \\
&+\beta_6 Orientation_i+\beta_7 Style_i+\beta_8 Duration_i
+\beta_9 Days_i+\varepsilon_i .
\end{aligned}
\end{equation}
\end{compactformula}
其中，$Y_i$ 表示第 $i$ 条视频的 $\log(\text{总互动量}+1)$。模型中的分类变量对应表 \ref{tab:variables-ai} 和表 \ref{tab:variables-content} 所列的 AI 使用特征和视频内容特征；$Duration_i$ 表示 $\log(\text{视频时长}+1)$，$Days_i$ 表示 $\log(\text{发布距采集日天数}+1)$，$\varepsilon_i$ 为随机误差项。

由于因变量为对数形式，分类变量系数可以通过 $\exp(\beta)-1$ 转换为相对于参照组的比例影响。连续变量系数则反映自变量对数变化与互动量对数变化之间的关系。

除主模型外，本文还估计包含省份固定效应的模型，用于考察地区差异是否影响主要结论。省份固定效应模型能够控制不同省份文旅账号基础、地方旅游资源和运营能力等相对稳定的地区差异，但由于会引入较多虚拟变量，本文将其作为稳健性补充，而不作为主模型。

\section{分析方法}

本文按照“描述特征--检验差异--解释影响--验证稳健性”的顺序进行数据分析。描述统计用于呈现点赞量、评论量、收藏量、转发量、总互动量、日均互动量、视频时长和发布距采集日天数的基本分布，同时分析各分类变量的频数和占比。

由于互动量分布右偏且不满足严格正态性，单因素差异检验采用 Kruskal-Wallis H 检验和 Mann-Whitney U 检验，用以判断各变量在表面上是否与传播效果存在差异关系。多元回归分析通过 OLS 模型同时纳入多个解释变量，识别各变量在控制其他因素后的净影响，是回答“每个指标是否有影响以及影响多少”的核心方法。

在高互动视频分析中，本文将总互动量位于样本前 25\% 的视频定义为高互动组，并采用组间比较和二分类 Logit 模型进行补充分析。组间比较用于呈现高互动样本的内容画像，二分类 Logit 则用于观察哪些变量更容易提高视频进入高互动组的可能性。该模型属于补充识别实验，主要服务于高互动样本解释，不替代以连续互动量为因变量的主回归模型。

稳健性检验从替换因变量、调整时间口径、处理极端值和加入省份固定效应四个角度展开。替换因变量时，分别使用点赞、评论、收藏和转发作为结果变量；调整时间口径时，引入日均互动量辅助判断采集时点差异是否改变结论；极端值处理则采用剔除总互动量 Top 1\%、剔除 Top 5\% 和 1\%-99\% 缩尾等方式重新估计模型；省份固定效应用于控制地区资源、账号基础和运营条件等相对稳定差异，检验核心内容变量是否仍保持基本解释方向。
